\documentclass[12pt]{article}
\usepackage[utf8]{inputenc}
\usepackage[T1]{fontenc}
\usepackage{hyperref}
\title{Central Pashto Language Resources for K-12 Curriculum Development}
\author{Compiled by Claude Sonnet 5.5 (claude-sonnet-5-5) for Andr\'as Kornai}
\date{}
\begin{document}
\maketitle

\section{Introduction}
Central Pashto (ISO 639-3 \texttt{pst}, Glottocode cent1973) is an Indo-European (Iranian) language of Pakistan, part of the Pashto macrolanguage, with about 6.5 million speakers. The search results did not target Central Pashto specifically; the items below are general Pashto resources, and their coverage of the Central variety is not verified. No Pashto-medium school curriculum materials were found in this search.

\section{Primary Educational Resources}
\begin{itemize}
\item Rahmon Inomkhojayev, \emph{Pashto: An Elementary Textbook}, two volumes with Pashto--English glossaries and Arabic-based alphabet instruction (aimed at adult learners of English-speaking background, not K-12). \url{https://www.booksamillion.com/p/Pashto/Rahmon-Inomkhojayev/9781589017733}
\end{itemize}

\section{Linguistic and Reference Materials}
\begin{itemize}
\item Anne David, \emph{Descriptive Grammar of Pashto and its Dialects} (De Gruyter, 2014). \url{https://www.degruyterbrill.com/document/doi/10.1515/9781614512318/html}
\item Awde and Sarwan, \emph{Pashto--English/English--Pashto Dictionary \& Phrasebook} (Hippocrene), about 5{,}000 romanized entries. \url{https://www.abebooks.com/9780781809726/Pashto-EnglishEnglish-Pashto-Dictionary-Phrasebook-Hippocrene-Phrasebooks-078180972X/plp}
\item Cambridge University Language Centre catalogue entry for a Pashto dictionary and phrasebook. \url{https://www.langcen.cam.ac.uk/resources/langp/data/pu_112.html}
\item Scribd list of Pashto dictionaries and grammars, including Bellew's 1867 historical dictionary and grammar (compiled list, quality unverified). \url{https://es.scribd.com/doc/288298711/Pashto-Dictionaries}
\end{itemize}

\section{Religious and Community Resources}
None found in this search.

\section{Digital and Audio Resources}
None found in this search.

\section{Curriculum Development Context}
The language is widely spoken, but the resources found are mostly commercial learner materials in English. Native-medium school textbooks, government curricula and audio would need to be sought through further searches (for example Pakistani provincial textbook boards), which this pass did not confirm.

\end{document}